\section{Perturbation Theory}

To compute expectation values of products of field operators, one can follow the Feynman
integral approach. The aim is to define this integral without trying to integrate over
field configurations, and in some cases this can be done rigorously by defining this special
integral as a linear operator on commuting (bosonic) or anticommuting (fermionic) variables
in such a way that the integral of polynomials respects Wick's rule. Bosonic field operators
can be represented with complex variables, this way their expectation values become complex
Gaussian integrals. Fermionic field operators on the other hand require a different
algebraic structure.

\subsection{Bosonic perturbation theory}

\subsection{Grassmann Algebras and Grassmann Integrals}

Consider the unital complex Grassmann algebra generated by a set of $2N$ anticommuting vectors:

\begin{equation}
	\begin{split}
	\label{def:grassmann}
		&\mathcal{G}_N = \mathrm{span}\{ 1, \psi_1, \bar\psi_1, \ldots, \psi_N, \bar\psi_N \} \\
		&\{ \psi_i, \psi_j \}         = 0 \quad \forall i,j=1,...,N \\
		&\{ \psi_i, \bar\psi_j \}     = 0 \quad \forall i,j=1,...,N \\
		&\{ \bar\psi_i, \bar\psi_j \} = 0 \quad \forall i,j=1,...,N
	\end{split}
\end{equation}

where $\{\psi, \phi\} := \psi\phi + \phi\psi$ is the anticommutator. The idea is to interpret
the elements of this algebra as (polynomials of) field operators, so we need to define integrals
and expectations of Grassmann variables, i.e. of elements of this algebra.

Given $\mathcal{G}_N$ the Grassmann algebra generated by
$\{\psi_1, ..., \psi_N, \bar\psi_1, ..., \bar\psi_N\}$, denote by
\begin{align*}
	\mathcal{G}_N^{-} &= \{
		\phi \in \mathcal{G}_N : [\phi, \psi] = 0 \;\forall \psi \in \mathcal{G}_N
	\} \\
		\mathcal{G}_N^{+} &= \{
		\phi \in \mathcal{G}_N : \{\phi, \psi\} = 0 \;\forall \psi \in \mathcal{G}_N
	\}
\end{align*}
the even (commuting) and odd (anticommuting) subsets of the algebra. It is easy to show that
these are subspaces (by linearity of commutators and anticommutators), and
\textcolor{red}{one can also show} that $\mathcal{G}_N$ can be decomposed as
\begin{align*}
	\mathcal{G}_N \cong \mathcal{G}_N^{+} \oplus \mathcal{G}_N^{-}.
\end{align*}
Odd elements satisfy a really important property. If $\psi \in \mathcal{G}_N^{+}$ it follows
that $\{\psi,\psi\} = 2\psi^2 = 0$. That is, all powers greater than one of an odd element vanish.

\begin{definition}
	Given $\mathcal{G}$ a Grassmann algebra, as defined above, and $f : \mathbb{C}^k \to \mathbb{C}$
	an analytic function given by its Taylor series centered at zero
	\begin{align*}
		f(\boldsymbol{z}) = \sum_{j=0}^\infty \sum_{|\alpha| = j}
		\frac{\partial^\alpha f(0)}{j!} \boldsymbol{z}^i
	\end{align*}
	we define the associated \textit{function of Grassmann variables} as
	\begin{align*}
		f(\boldsymbol{\psi}) = \sum_{j=0}^\infty \sum_{|\alpha| = j}
		\frac{\partial^\alpha f(0)}{j!} \boldsymbol{\psi}^i
	\end{align*}
	where $\boldsymbol{\psi} \in \mathcal{G}_N \times ... \times \mathcal{G}_N$ is the list of
	$k$ Grassmann variables.
\end{definition}

\begin{remark}
	If $\psi \in \mathcal{G}_N^{+}$ it follows that any function of Grassmann variables evaluated
	in $\psi$ has at most first order terms, as all higher orders will always contain a factor
	of $\psi^2 = 0$. The most important example to us is the exponential of an odd element:
	\begin{align*}
		e^{\psi} = 1 + \psi + \frac{1}{2}\psi^2 + ... = 1 + \psi
	\end{align*}
\end{remark}

\begin{remark}
	We could define a Grassmann algebra more "directly" as the ring of polynomials in
	anticommuting variables over a field, and choose $\mathbb{C}$ as base field to get exactly
	the structure we need. Indeed, the Grassmann algebra can be viewed as the free algebra
	generated by the variables corresponding to the modes of the physical system, quotiented
	by the ideal generated by anticommutators, i.e.
	\begin{align*}
		\mathcal{G}_N \coloneq
		\frac{\mathbb{C} \langle \psi_1, ..., \psi_N, \bar\psi_1, ..., \bar\psi_N \rangle}
		{ \mathrm{span}\{ \psi_i\psi_j + \psi_j\psi_i,\, \psi_i\bar\psi_j + \bar\psi_j\psi_i,\, \bar\psi_i\bar\psi_j + \bar\psi_j\bar\psi_i : i,j = 1, ..., N \} }
	\end{align*}
\end{remark}

We now introduce integrals of functions of Grassmann variables. Note that the definition is
going to be purely algebraic and, in particular, motivated by our need to reproduce specific
computational rules. Indeed, we'll see that we can define expectations of functions of
Grassmann variables in a way that reproduces exactly Wick's rule for expectations of fermionic
field operators.

We ask, first of all, for the integral to be a linear operator on a Grassmann algebra.
Then, specifying the image of a basis for the algebra will uniquely determine the
operator.
\begin{definition}
	Given $\{1, \psi_1, ..., \psi_N, \bar\psi_1, ..., \bar\psi_N\}$ the basis for the Grassmann algebra,
	as above, the \textit{Grassmann integral} is the linear operator that acts on this basis as
	\begin{align*}
		&\int\dx\psi_i\,\psi_j = \int\dx\bar\psi_i\,\bar\psi_j \coloneq \delta_{ij}, \\
		&\int\dx\psi_i = \int\dx\bar\psi_i \coloneq 0.
	\end{align*}
	Repeated integration is denoted by a single integral symbol and multiple "differentials".
\end{definition}

\begin{example}
	In $\mathcal{G}_2$, consider the exponential function. Using the definition, we can
	evaluate the following integrals to see how this operator behaves.
	\begin{itemize}
		\item $\int\dx\psi_1\,e^{\psi_1\bar\psi_2} = \int\dx\psi_1\,(1 + \psi_1\bar\psi_2) = \bar\psi_2$;
		\item $\int\dx\bar\psi_2\,e^{\psi_1\bar\psi_2} = \int\dx\bar\psi_2\,(1 - \bar\psi_2\psi_1) = -\psi_1$;
		\item $\int\dx\psi_1\dx\bar\psi_2\,e^{\psi_1\bar\psi_2} = -1$.
	\end{itemize}
\end{example}

\begin{remark}
	Repeated integration, i.e. the composition of multiple integral operators, inherits some form of
	anticommutativity when acting on the basis elements. Indeed in $\mathcal{G}_1$, for example, we have
	\begin{align*}
		\int\dx\psi\dx\bar\psi\,\psi\bar\psi &= -\int\dx\psi\dx\bar\psi\,\bar\psi\psi = -1 \\
											 &= -\int\dx\bar\psi\dx\psi\,\psi\bar\psi.
	\end{align*}
	And given the definition of Grassmann functions, this property will be quite useful to us.
\end{remark}

Having defined integration, we can now introduce a \textit{Gaussian measure} (in a purely algebraic
sense). To make the notation ligher, we define a "dot product", without any geometric meaning, as
$\boldsymbol{\psi}\cdot\boldsymbol{\phi} \coloneq \sum_{i = 1}^k \psi_i\phi_i$.

\begin{definition}
	$G \in \mathrm{Mat}(N \times N, \mathbb{C})$ is said to be a \textit{Grassmann covariance}
	if $G \in \mathrm{GL}(N,\mathbb{C})$.
\end{definition}

\begin{lemma}
	Let $G$ be a Grassmann covariance, $\boldsymbol{\psi} = (\psi_1, \ldots, \psi_N)$ and
	$\boldsymbol{\bar\psi} = (\bar\psi_1, \ldots, \bar\psi_N)$. Then
	\begin{align*}
		\int\left( \prod_{i=1}^N \dx\bar\psi_i \dx\psi_i \right)
		e^{-\boldsymbol{\bar\psi} \cdot G^{-1} \boldsymbol{\psi}} = \mathrm{det}(G^{-1})
	\end{align*}
\end{lemma}
\begin{proof}
	We carry out the calculation explicitly.
	\begin{align*}
		\int\left( \prod_{i=1}^N \dx\bar\psi_i \dx\psi_i \right)
		e^{-\boldsymbol{\bar\psi} \cdot G^{-1} \boldsymbol{\psi}}
		&= \int\left( \prod_{i=1}^N \dx\bar\psi_i \dx\psi_i \right)
		\frac{(-1)^N}{N!}\left( \sum_{i,j=1}^N \bar\psi_i (G^{-1})_{ij} \psi_j \right)^N \\
		&= \int\left( \prod_{i=1}^N \dx\bar\psi_i \dx\psi_i \right)
		\frac{1}{N!}\left( \sum_{i,j=1}^N \bar\psi_i (G^{-1})_{ij} \psi_j \right)^N \\
		&= \int\left( \prod_{i=1}^N \dx\bar\psi_i \dx\psi_i \right)
		\frac{1}{N!} \sum_{\sigma,\tau \in \mathrm{S}_N} \prod_{j=1}^N (G^{-1})_{\sigma(j),\tau(j)} \bar\psi_{\sigma(j)} \psi_{\tau(j)}
		= \star
	\end{align*}
	Where $\mathrm{S}_N$ is the set of permutations of N elements. We can split the product
	into a product of matrix elements times a product of Grassmann variables, then we use that
	\begin{align*}
		\prod_{j=1}^N \bar\psi_{\sigma(j)}\psi_{\tau(j)} = (-1)^{\mathrm{sgn}(\sigma) + \mathrm{sgn}(\tau)} \prod_{j=1}^N \bar\psi_j \psi_j
	\end{align*}
	to find that
	\begin{align*}
		\star &= \frac{1}{N!} \sum_{\sigma,\tau \in \mathrm{S}_N} (-1)^{\mathrm{sgn}(\sigma) + \mathrm{sgn}(\tau)}
		\prod_{j=1}^N (G^{-1})_{\sigma(j),\tau(j)} \\
		&= \frac{1}{N!} \sum_{\sigma \in \mathrm{S}_N} (-1)^{\mathrm{sgn}(\sigma)}
		\prod_{j=1}^N (G^{-1})_{j,\sigma{j}} \\
		&= \mathrm{det}(G^{-1})
	\end{align*}
\end{proof}

\begin{definition}
	Let $G$ be a Gaussian covariance, and $\{ \psi_1, \ldots, \psi_N, \bar\psi_1, \ldots, \bar\psi_N\}$
	the usual Grassmann variables. A \textit{Grassmann Gaussian measure} is the Grassmann
	integral measure defined by
	\begin{align*}
		P_G(\dx\psi) \coloneq \frac{1}{\mathrm{det}(G^{-1})}
		\left( \prod_{j=1}^N \dx\bar\psi_j \dx\psi_j \right)
		e^{-\boldsymbol{\bar\psi} \cdot G^{-1} \boldsymbol{\psi}}
	\end{align*}
\end{definition}

\begin{remark}
	The factor in front of the differentials guarantees a "normalized" measure, i.e.\\
	$\int P_G(\dx\psi) = (1/\mathrm{det}(G^{-1})) \mathrm{det}(G^{-1}) = 1$.
\end{remark}

We will prove the following statement later.

\begin{lemma}
	With the same notation used so far,
	\begin{align*}
		\int P_G(\dx\psi) \psi_i \bar\psi_j = G_{ij}
	\end{align*}
\end{lemma}

\begin{lemma}
	Let $D \in \mathrm{GL(2,\mathbb{C})}$, $\boldsymbol{\varphi}$ and $\boldsymbol{\eta}$
	two vectors of Grassmann variables. Then, under the following change of variables
	\begin{align*}
		\begin{pmatrix} \boldsymbol{\psi} \\ \boldsymbol{\bar\psi} \end{pmatrix}
		= D \begin{pmatrix} \boldsymbol{\varphi} \\ \boldsymbol{\bar\varphi} \end{pmatrix}
		+ \begin{pmatrix} \boldsymbol{\eta} \\ \boldsymbol{\bar\eta} \end{pmatrix}
	\end{align*}
	the integral of a function of Grassmann variables $f$ changes as follows
	\begin{align*}
		\int \left( \prod_{i} \dx\bar\psi_i \psi_i \right) f(\boldsymbol{\psi},\boldsymbol{\bar\psi})
	= \mathrm{det}(G^{-1}) \int \left( \prod_{i} \dx\bar\varphi_i \varphi_i \right) f(D(\boldsymbol{\varphi},\boldsymbol{\bar\varphi})+(\boldsymbol{\eta},\boldsymbol{\bar\eta}))
	\end{align*}
\end{lemma}
\begin{proof}
	\textcolor{red}{check proof in the notes (D is defined only in the diagonal NxN
	blocks? How does he calculate "differentials" of the new variables?)}.
\end{proof}

\begin{lemma}
	Let $G \in \mathrm{GL}(N,\mathbb{C})$, then
	\begin{align*}
		\int P_G(\dx\psi)\,e^{\boldsymbol{\bar\eta \cdot \psi} + \boldsymbol{\bar\psi \cdot \eta}}
			= e^{\boldsymbol{\bar\eta \cdot} G \boldsymbol{\eta}}
	\end{align*}
\end{lemma}
\begin{proof}
	\textcolor{red}{it was an exercise.}
\end{proof}

Before defining expectations, we need one more operator on a Grassmann algebra to mimic the
behaviour of these same integrals done in the complex variables case. That is, we need to introduce
derivatives of Grassmann functions. If we had a polynomial of degree N in N variables, then
deriving this polynomial in all of the variables should return the maximal degree term. Then
this becomes a natural choice for the definition of Grassmann derivatives, as we defined
functions to be nothing more than Taylor series, and evaluating a function in the elements of
the basis gets us exactly in this situation. But recall that integrating a Grassmann function
in all variables (i.e., in all the basis elements) returns exactly the same thing, the coefficient
of the maximal degree term. Therefore, this natural definition for Grassmann derivatives
reads as follows.

\begin{definition}
	Let $f(\boldsymbol{\psi},\boldsymbol{\bar\psi})$ be a Grassmann function, the
	\textit{Grassmann derivatives} of $f$ w.r.t. the basis elements are the linear operators
	defined as
	\begin{align*}
		\frac{\partial}{\partial\psi_i}f \coloneq \int\dx\psi_i\,f(\boldsymbol{\psi},\boldsymbol{\bar\psi})
	\end{align*}
\end{definition}

\begin{proof}[Proof (Lemma 2.2)]
Recalling the "anticommutativity property" of Grassmann integrals, it is clear that the same
property holds for derivatives too. We can make use of Grassmann derivatives to prove this
lemma in a way that is closer to the complex variables case, and this motivates the introduction
of Grassmann derivatives. Note that, given a vector of Grassmann variables $\boldsymbol{\eta}$,
\begin{align*}
	\psi_i\bar\psi_j e^{-\boldsymbol{\bar\psi}\cdot G^{-1} \boldsymbol{\psi}}
		= -\frac{\partial}{\partial\bar\eta_i}\frac{\partial}{\partial\eta_j}
		e^{-\boldsymbol{\bar\psi}\cdot G^{-1} \boldsymbol{\psi}
		+ \boldsymbol{\bar\eta} \cdot \boldsymbol{\psi} + \boldsymbol{\bar\psi}\cdot\boldsymbol{\eta}}
		\bigg|_{\boldsymbol{\eta} = \boldsymbol{\bar\eta} = 0}.
\end{align*}
We can substitute this in the integral to get that
\begin{align*}
	\int P_G(\dx\psi)\, \psi_i\bar\psi_j
	&= -\frac{\partial}{\partial\bar\eta_i}\frac{\partial}{\partial\eta_j} \int P_G(\dx\psi)
	e^{\boldsymbol{\bar\eta} \cdot \boldsymbol{\psi} + \boldsymbol{\bar\psi}\cdot\boldsymbol{\eta}}
	\bigg|_0 \\
	&= -\frac{\partial}{\partial\bar\eta_i}\frac{\partial}{\partial\eta_j}
	e^{\boldsymbol{\bar\eta} \cdot G \boldsymbol{\eta}} \bigg|_0 \\
	&= -\frac{\partial}{\partial\bar\eta_i}\frac{\partial}{\partial\eta_j}
	(\bar\eta_i G_{ij} \eta_j) \bigg|_0 \\
	&= G_{ij}
\end{align*}
which is exactly the thesis.
\end{proof}

With this useful lemma proven, we get to the most important theorem of this section: Wick's
theorem. It is this result that allows us to actually interpret Grassmann Gaussian integrals
as "expectations of fermionic field operators", because it shows that these integrals exactly
reproduce the combinatorics of fermionic expectations.

\begin{theorem}[Wick's rule]
	Let $M \in \mathbb{N}$ be the degree of a monomial in Grassmann variables, and let
	$I = \{a_1, \ldots, a_M\}$ and $J = \{b_1, \ldots, b_M\}$ be two sets of M indices. The
	integral of a monomial is
	\begin{align*}
		\int P_G(\dx\psi) \left( \prod_{\alpha \in I, \beta \in J} \psi_{\alpha}\bar\psi_{\beta} \right)
		= \sum_{\sigma \in \mathrm{S}_N} (-1)^{\mathrm{sgn}(\sigma)} \prod_{j=1}^M G_{a_j,b_{\sigma(j)}}
	\end{align*}
\end{theorem}
\begin{proof}
	\textcolor{red}{check proof of Isserlis' theorem, and check the way the product on the LHS is written.}
\end{proof}

\textcolor{red}{I REWORKED THE STRUCTURE OF THIS SECOND CHAPTER!}.
It follows that calling a Grassmann Gaussian integral a "fermionic expectation" makes sense.
Our task now is to give important results that greatly simplify the calculation of such
expectations. In particular, the aim is to represent these calculations using Feynman diagrams,
which are a tool to perform the combinatorics of this exact type of problem, and because these
integrals are supposed to be a representation of expectations of physical quantities, Feynman
diagrams will have a clear physical interpretation that we will discuss in due time.

The results we are about to give concern both bosonic and fermionic systems, and the differences
often lie only in sign changes between expressions. In fact, the rules behind Feynman diagrams
are going to be the same, what will change however is the precise formula to evaluate such
diagrams in a complete calculation.

\subsection{Fermionic perturbation theory}

\subsection{Feynman diagrams}
